%% GENERATED by python/paper/build.py from results/ -- do not edit by hand.
%% Rebuild:  .venv\Scripts\python.exe python\paper\build.py --only US_Ageing
%% Source:   results/shocks/US_shocks.csv
\begin{table}[!htb]
\centering
\begin{threeparttable}
\caption{The effect of ageing in US -- 2020}
\label{table:US:ageing}
\renewcommand{\arraystretch}{1.25}
\begin{tabularx}{.9\textwidth}{p{3cm}YYY}
\toprule
\textbf{Scenario} & \textbf{Tax rate} & \textbf{Savings rate} & \textbf{Avg. workweek} \\
\midrule 
Baseline & 14.43\% & 21.96\% & 39.39 \\[1.25ex]
\multicolumn{4}{l}{\textit{Full effect:}} \\\hline
Mild ageing\tnote{a} & 18.45\% & 20.54\% & 39.37 \\
Acute ageing\tnote{b} & 23.11\% & 18.99\% & 39.30 \\[1.25ex]
\multicolumn{4}{l}{\textit{Economic equilibrium effect:}} \\\hline
Mild ageing\tnote{a} & 14.43\% & 21.96\% & 39.84 \\
Acute ageing\tnote{b} & 14.43\% & 21.96\% & 40.36 \\[1.25ex]
\bottomrule
\end{tabularx}
\begin{tablenotes}
\footnotesize
\item[a] The ``mild ageing'' scenario refers to the case with $\nu_t$ set at $(1+\nu_t^{base})/2$ from 2020 and onward.
\item[b] The ``acute ageing'' scenario refers to $\nu_t = 1$ from 2020 and onward.
\end{tablenotes}
\end{threeparttable}
\end{table}
